\documentclass[11pt,a4paper]{article}
\usepackage{amsmath,amssymb,amsthm}
\usepackage[margin=2.5cm]{geometry}
\usepackage{hyperref}

\newtheorem{definition}{Definition}[section]
\newtheorem{theorem}[definition]{Theorem}
\newtheorem{proposition}[definition]{Proposition}
\newtheorem{lemma}[definition]{Lemma}
\newtheorem{corollary}[definition]{Corollary}
\newtheorem{conjecture}[definition]{Conjecture}
\newtheorem{remark}[definition]{Remark}

\title{Hyperseed Formalization:\\
  Toward a Truly Decentralized Digital Provenance Layer}
\author{ProtomegaTron (automated formalization)}
\date{2026-09-17}

\begin{document}
\maketitle

\begin{abstract}
We formalize the intellectual structure of Ben Goertzel's Substack article
``Toward a Truly Decentralized Digital Provenance Layer'' (2026-08-25) within
the Hyperseed ontology. The article introduces OpenWater, an open,
decentralized media provenance framework that shifts from AI-based deepfake
detection (a losing arms race) to cryptographic provenance (signed claims
assembled into a sheaf of evidence). Each design principle maps onto
Hyperseed structures: signed claims as local sections, trust policy as
per-observer fiber evaluation, decentralized resolution as distributed
fiber bundle, the liar's dividend as provenance vacuum exploitation, and
design refusals as anti-cocycle-collapse guards.
\end{abstract}

\tableofcontents

%% =================================================================
\section{Detection vs.\ provenance: observables vs.\ directed history}
\label{sec:detection}
%% =================================================================

\begin{proposition}[Detection as losing arms race --- claims 1--3]
\label{prop:arms-race}
AI-based deepfake detection is a co-evolutionary arms race where the
faker side has structural advantages:
\begin{enumerate}
\item A detector returns a probability, not a history --- it says how
  likely the current state is to be AI-generated, without reconstructing
  the path that produced it.
\item New generators learn to evade old detectors; compression and
  re-editing confuse classifiers; a capable attacker trains against the
  detector itself.
\item The detection approach lives at the 0-cell level of the directed
  type: it examines observable properties of the current state without
  accessing higher cells.
\end{enumerate}
\end{proposition}

\begin{definition}[Provenance as directed history --- claim 25]
\label{def:provenance}
\emph{Provenance} asks about the directed history of an artifact: what
path did it take through the directed type of transformations?
\[
  \text{Provenance}(a) = \mathcal{H}_a = (h_0 \xrightarrow{t_1} h_1
  \xrightarrow{t_2} \cdots \xrightarrow{t_n} h_n)
\]
where each $h_i$ is a state and each $t_i$ is a signed transformation
claim. The shift from detection to provenance is the shift from
examining 0-cells (current observable state) to examining 1-cells
(transformation history) and higher cells (justifications, cross-references).
\end{definition}

%% =================================================================
\section{The liar's dividend as provenance vacuum exploitation}
\label{sec:liar}
%% =================================================================

\begin{proposition}[Provenance vacuum --- claims 4, 27]
\label{prop:vacuum}
The \emph{liar's dividend}: once nothing can be authenticated, the most
powerful actor in any dispute can deny whatever evidence is inconvenient.
In Hyperseed terms, when the base space $B$ has no provenance sheaf
$\mathcal{E}$ over it:
\[
  \mathcal{E} = \emptyset \;\Rightarrow\;
  \text{any actor can assert any section } s : U \to E
  \text{ without contradiction}.
\]
The second-order effect (real material dismissed as fake) may be more
corrosive than the first (fake presented as real), because the absence
of provenance structure means there is no descent data to arbitrate
between competing claims.
\end{proposition}

%% =================================================================
\section{OpenWater as provenance sheaf}
\label{sec:openwater}
%% =================================================================

\begin{definition}[Provenance sheaf --- claims 5--9, 23]
\label{def:sheaf}
The OpenWater framework constructs a provenance sheaf $\mathcal{E}$
over the space of artifacts:
\begin{enumerate}
\item \textbf{Signed claims as local sections}: each producer or
  transformer (camera, AI model, editor, publisher, agent) signs
  precise claims about what it did, producing a local section
  $e_i : U_i \to \text{Evidence}$.
\item \textbf{Interoperable credential}: local sections are packaged
  into C2PA-compatible credentials --- the shared format that enables
  transition functions between different signers' claims.
\item \textbf{Invisible watermark}: a robust invisible watermark
  provides the persistent link from artifact to credential, surviving
  metadata stripping --- the watermark is the map back to the sheaf
  even when the explicit coordinate system (metadata) is destroyed.
\item \textbf{Decentralized resolution}: credentials stored by many
  independent services --- a distributed fiber bundle, not a single
  mandatory database.
\item \textbf{Trust policy}: each observer applies their own trust
  policy to weight the evidence --- per-observer fiber evaluation,
  not a global truth function.
\end{enumerate}
\end{definition}

\begin{theorem}[Shared grammar, not shared verdict --- claim 20]
\label{thm:grammar}
The OpenWater design produces a shared grammar of evidence: all
participants speak the same protocol (C2PA-compatible signed claims),
but different observers project the evidence fiber onto different
evaluation subspaces according to their trust policy. A science journal,
indigenous media network, national archive, dissident collective, and
social platform can accept different sets of signers while speaking the
same underlying protocol. This is the sheaf-theoretic analog of a
universal covering space with different group actions.
\end{theorem}

%% =================================================================
\section{Trust policy as per-observer fiber evaluation}
\label{sec:trust}
%% =================================================================

\begin{definition}[Per-observer evaluation --- claims 9, 19, 24, 29]
\label{def:observer}
Provenance verdicts are fiber-valued, not binary. For an artifact $a$
with evidence fiber $E_a = (e_1, e_2, \ldots, e_k)$ spanning $k$
evidence dimensions:
\[
  \text{Verdict}_O(a) = \pi_O(E_a) = \sum_{i=1}^{k} w_i^{(O)} \cdot e_i
\]
where observer $O$ applies weights $w_i^{(O)}$ from their trust bundle.
Different observers produce different projections of the same fiber ---
no scalar collapse to a single truth/false value. This is the same
anti-scalar-collapse stance as:
\begin{itemize}
\item The $(f,c)$-lossy critique (visible STV fusion vs.\ per-packet
  provenance),
\item Rejection of a single ``personhood score'' in the AI Rights
  article,
\item Scenario-over-doom-probability methodology.
\end{itemize}
\end{definition}

%% =================================================================
\section{Design refusals as anti-cocycle-collapse guards}
\label{sec:refusals}
%% =================================================================

\begin{proposition}[Four refusals --- claims 10--14, 30]
\label{prop:refusals}
OpenWater's four design refusals each prevent a specific mode of
cocycle collapse:
\begin{enumerate}
\item \textbf{No global truth authority}: prevents collapse of
  the distributed sheaf to a single authoritative section
  (centralized cocycle).
\item \textbf{No unstaged-event claim}: prevents collapse of the
  provenance question (``who signed what?'') to a content question
  (``did this really happen?'') --- keeping the sheaf about
  transformation history, not about ground truth.
\item \textbf{No civil identity requirement}: prevents collapse
  of the participant fiber to a single identity dimension ---
  preserving pseudonymous and anonymous contributions.
\item \textbf{No absence-as-guilt}: prevents collapse of the
  complement (unmarked artifacts) into a negative category ---
  maintaining that the absence of provenance data is absence of
  evidence, not evidence of absence.
\end{enumerate}
These refusals are the design discipline that keeps a provenance layer
from mutating into censorship infrastructure.
\end{proposition}

%% =================================================================
\section{AI for reputation as holonomy detection}
\label{sec:ai-reputation}
%% =================================================================

\begin{proposition}[Reputation holonomy --- claims 15--16, 28]
\label{prop:reputation}
The core AI role in OpenWater is reputation management: detecting
sophisticated reputation-gaming by participants. In Hyperseed terms,
this is holonomy detection in the trust bundle:
\[
  \text{Hol}(\gamma_{\text{gaming}}) \neq \text{id}
\]
where $\gamma_{\text{gaming}}$ is a path through the reputation space
where an actor fakes good behavior (accumulates undeserved reputation)
and then exploits it. The nontrivial holonomy --- the actor returns to
``trustworthy'' status via a deceptive path that doesn't preserve the
trust fiber's value --- is exactly the pattern AI is trained to detect.

The role is specifically \emph{not} deepfake detection (which is the
losing arms race). AI polices the reputation layer; cryptography and
provenance handle the evidence layer.
\end{proposition}

%% =================================================================
\section{C2PA compatibility: format vs.\ infrastructure}
\label{sec:c2pa}
%% =================================================================

\begin{proposition}[Necessary not sufficient --- claims 17--18]
\label{prop:c2pa}
An open format (C2PA) is necessary but not sufficient. If repositories,
watermark resolution, key histories, revocation lists, and trust
decisions are centralized, the system rebuilds the same bottleneck
one layer up --- a centralized infrastructure wearing an open-standard
mask. All layers must be plural:
\[
  \text{Open format} + \text{centralized infra} \cong \text{centralized system}
\]
\[
  \text{Open format} + \text{decentralized infra} \cong \text{distributed fiber bundle}
\]
\end{proposition}

%% =================================================================
\section{Cross-references}
%% =================================================================

\begin{remark}[Connections to other formalizations]
\begin{itemize}
\item \textbf{Proof of Humanity} (2026-08-26): the companion article.
  OpenWater provides the provenance infrastructure; proof-of-humanity
  is $E_\pi$ restricted to the identity fiber. Same composite-proof
  sheaf architecture.
\item \textbf{Sanders Ban} (2026-09-05): selection effect as provenance
  destruction. The liar's dividend is the passive version of what the
  Sanders ban does actively.
\item \textbf{AI Rights} (2026-09-02): evaluation ecology as sheaf of
  assessment sections. The same multi-observer, anti-scalar-collapse
  architecture applied to personhood evaluation.
\item \textbf{Seven Flavors} (2026-07-01): burning observability as
  provenance destruction. Building uninterpretable models is the AI
  development version of stripping metadata.
\item \textbf{Coxon/EA} (2026-09-13): auditor capture as cocycle
  failure. OpenWater's decentralized resolution prevents the single
  auditor capturing the entire system.
\item \textbf{$(f,c)$-lossy critique} (LTM 5a4d150b): per-observer
  fiber evaluation as anti-scalar-collapse.
\end{itemize}
\end{remark}

\begin{conjecture}[Watermark robustness as fiber persistence]
The invisible watermark's robustness under compression, cropping, and
re-encoding is a fiber persistence property: the link from artifact
to provenance sheaf must survive transformations that destroy explicit
metadata. Conjecture: the minimum watermark capacity needed for sheaf
recovery is proportional to the number of independent evidence dimensions
$k$ in the provenance fiber, not to the total amount of evidence ---
because the watermark only needs to point to the sheaf entry, not to
carry the evidence itself.
\end{conjecture}

\end{document}
